% sec-06-funding-proposals.tex - Section 6, Funding Proposals.  A main section.
% FULL stage.  Figures 17, 18, 19 and 20 drawn; Tables 17 to 20 populated from
% funding/capitalization-plan/final-capital/publication,
% funding/pdac-funding-applications/final-apply/publication, and
% funding/RFA-RM-27-001-v2.  The 36,330 dollar valuation is reported as
% projected, never as estimated, in the prose and in every row of Table 20.
\section{Funding Proposals}
\label{sec:funding}

\subsection{Ten applications, one method}

In August 2026 ten funding applications were filed for a Phase 1 trial of an
LLM-advised robotic pancreaticoduodenectomy with perioperative daraxonrasib,
written throughout as an independent scientist \cite{tenapps2026}. Nine of the
ten address mechanisms that fund a person or an institution. One, application
05, addresses the only mechanism in the set built for a company, and it was the
shortest of the ten. Each application answers a specific clause: application 01
cites the Director's Pioneer Award as the model for ``scaling long-horizon
grants for the best and brightest'', application 03 addresses ``the first
federal program designed to fund independent research organizations outside
traditional academia'', and application 05 answers the sentence that
``programs like SBIR open doors for technician-founded ventures''
\cite{tenapps2026, kratsios2026sciencegoldenage}. Table 17 lists all ten with
the mechanism class and the ask each carries.

\begin{apptable}
\begin{tabularx}{\textwidth}{>{\raggedright\arraybackslash}p{0.85cm} >{\raggedright\arraybackslash}p{4.15cm} >{\raggedright\arraybackslash}p{2.75cm} >{\raggedright\arraybackslash}p{1.95cm} >{\raggedright\arraybackslash}X}
\headrow \thead{No.} & \thead{Recipient} & \thead{Mechanism class} & \thead{Term} & \thead{Ask}\\
\midrule
01 & NIH Common Fund, Pioneer Award & Person-based & 5 years & \$3.5M \\
02 & ARPA-H mission office & Organization & 36 months & \$2.1M \\
03 & NSF TIP Directorate, X-Labs & Organization & 5 years & \$3.5M \\
04 & DOE Office of Science, Genesis Mission & Partnership & 5 years & \$3.5M \\
05 & NIH SEED, SBIR and STTR & Small business & 9 plus 24 months & \$1.6M \\
06 & Foundation for the NIH, AMP & Partnership & 5 years & \$3.5M \\
07 & HHMI Investigator Program & Person-based & 7 years & \$0.7M per year \\
08 & NCI CTEP & Concept review & 5 years & \$3.5M \\
09 & Convergent Research, FRO program & Organization & 5 years & \$3.5M \\
10 & UC San Diego Moores Cancer Center & Partnership & One meeting & None \\
\end{tabularx}
\end{apptable}
\vspace{-0.6cm}
\tabcap{Table 17. The ten applications filed in August 2026, the mechanism class\\each
addresses, the term requested, and the ask carried on each.}

Figure 17 places all fourteen 2026 funding artifacts on one calendar, and
Figure 20 draws the machinery that produced them. Four of the fourteen closed
within one week of the work that fed them, which is the production-rate claim in
its most checkable form.

% ---------------------------------------------------------------------------
% FIGURE 17.  mermaid-type, gantt.  One linear 74-day axis, production bar and
% darker deposit bar abutting on each row, four ringed deposits.
% ---------------------------------------------------------------------------
\begin{appfloat}[!tb]
\begin{appfig}[mermaid-type / gantt]
\ptitle{0}{0.85}{Funding artifact production and deposit, June to August 2026}
\axisx{0}{12.60}{0.30}
\foreach \x/\l in {0/{Jun 15},2.52/{Jun 30},5.04/{Jul 15},7.56/{Jul 30},10.08/{Aug 14},12.60/{Aug 29}}{
  \draw[trialslate,line width=0.3pt] (\x,0.30) -- (\x,0.40);
  \node[font=\tiny\sffamily,anchor=south] at (\x,0.42) {\l};}
\node[mmlanetitle,anchor=east] at (-0.15,-0.44) {Enabling};
\node[mmlanetitle,anchor=east] at (-0.15,-2.48) {RFA};
\node[mmlanetitle,anchor=east] at (-0.15,-3.50) {Apply};
\node[mmlanetitle,anchor=east] at (-0.15,-4.18) {Capital};
\foreach \i/\lab/\p0/\p1/\d1 in {%
  0/{Phase 1 protocol}/0.67/1.00/1.17,
  1/{Phase 2 protocol}/1.17/1.50/1.67,
  2/{PI adoption guide}/1.50/1.84/2.01,
  3/{Document guidance paper}/2.01/2.35/2.52,
  4/{Phase 1 IND}/2.18/2.85/3.02,
  5/{RFA v1 production}/3.19/3.69/3.86,
  6/{RFA v2 production}/4.03/4.53/4.70,
  7/{Patient robot advocacy}/7.39/7.89/8.06,
  8/{Ten application set}/8.73/10.07/10.41,
  9/{Capitalization plan}/9.90/10.40/10.57}{
  \node[anchor=east,font=\tiny\sffamily] at (-1.55,{-0.30-0.34*\i}) {\lab};
  \fill[trialburgp] (\p0,{-0.30-0.34*\i-0.11}) rectangle (\p1,{-0.30-0.34*\i+0.11});
  \draw[trialburg,line width=0.3pt] (\p0,{-0.30-0.34*\i-0.11}) rectangle (\p1,{-0.30-0.34*\i+0.11});
  \fill[trialburgl] (\p1,{-0.30-0.34*\i-0.11}) rectangle (\d1,{-0.30-0.34*\i+0.11});
  \draw[trialburg,line width=0.3pt] (\p1,{-0.30-0.34*\i-0.11}) rectangle (\d1,{-0.30-0.34*\i+0.11});}
% today marker
\draw[trialchar,dashed,line width=0.5pt] (12.10,0.30) -- (12.10,-3.40);
\node[font=\tiny\sffamily,anchor=south,text=trialchar] at (12.10,0.32) {Aug 14, 2026};
\node[mmmid,text width=32mm] at (2.10,-4.35) {Fourteen artifacts in 74 days};
\legkey{6.20}{-4.35}{trialburgp}{production window}
\legkey{9.10}{-4.35}{trialburgl}{deposit}
\pnote{-1.55}{-5.05}{Production and deposit bars share a row and abut, so each artifact reads as one
object whose right end is its\\deposit date. The deposit segment is always the
darker fill, which puts the eye on the date a funder reads.}
\end{appfig}
\vspace{-0.6cm}
\figcaption{Figure 17. Fourteen funding artifacts deposited between June and
August 2026,\\their production windows, and the four deposits that closed within
one week.}
\end{appfloat}

\subsection{The capitalization plan}

The capitalization plan of August 11, 2026 converts the author from independent
scientist to small-business operator and rewrites application 05 as the document
it should have been \cite{capplan2026}. It states its arithmetic without
softening: ``The identical four-layer budget costs \$3,500,000 of direct work
over sixty months and \$1,396,000 of direct work over thirty-three. The award
that carries the second is \$1,606,000, a 15.0 percent overhead and fee load
against an SBIR allowance of 40 percent that the company does not claim. Routed
through a university at a 57 percent facilities and administrative rate, the
same direct work would cost a funder \$2,137,000'' \cite{capplan2026}. And it
states the bridge: ``Federal money of \$1,606,000 sits below a firewall;
contributed drug, operating room and pathology sit unpriced in the middle; and
\$5,900,000 of private capital sits above, raised only after a milestone closes.
That is 3.67 to one on cash alone'' \cite{capplan2026, kratsiosvought2026fy2028priorities}.
Table 18 collects the ten quantities the plan turns on; Figure 18 prices the
same four-layer budget three ways against a fourth column that routes the
identical direct work through an institution, and Figure 19 shows where each
port of the award lands and which two work packages it does not reach.

\begin{apptable}
\begin{tabularx}{\textwidth}{>{\raggedright\arraybackslash}p{5.55cm} >{\raggedright\arraybackslash}p{2.55cm} >{\raggedright\arraybackslash}X}
\headrow \thead{Quantity} & \thead{Value} & \thead{What it governs}\\
\midrule
SBIR Phase I & \$306,000 & 9 months, 5 milestones \\
SBIR Phase II & \$1,300,000 & 24 months, 7 milestones \\
Total award & \$1,606,000 & 33 months \\
Direct work inside the award & \$1,396,000 & The work the award pays for \\
Direct work outside the award & \$2,104,000 & The rest of the five-year program \\
Five-year direct program & \$3,500,000 & 60 months, direct only \\
Company overhead and fee load & 15.0 percent & Against a 40 percent allowance not claimed \\
Same direct work, university route & \$2,137,000 & At a 57 percent facilities and administrative rate \\
Private capital above the firewall & \$5,900,000 & 3.67 to one on cash alone \\
Staffing at full Phase II & 2.6 FTE & The team the plan proposes \\
\end{tabularx}
\end{apptable}
\vspace{-0.6cm}
\tabcap{Table 18. The ten capitalization quantities the plan turns on, and what\\each
one governs in the proposed thirty-three month program.}

% ---------------------------------------------------------------------------
% FIGURE 18.  d2-type, grid.  One row height, five fixed column widths, no
% edges, cut once by a charcoal rule between scope rows and money rows.
% ---------------------------------------------------------------------------
\begin{appfloat}[!tb]
\begin{appfig}[d2-type / grid]
\def\rh{0.66}
\node[d2cellh,minimum width=3.8cm,minimum height=\rh cm,anchor=north west] at (0,0) {Cost layer};
\foreach \i/\h in {0/{SBIR Phase I, 9 months},1/{SBIR Phase II, 24 months},
  2/{Five year direct program},3/{Same work, university route}}{
  \node[d2cellh,minimum width=2.75cm,minimum height=\rh cm,anchor=north west]
    at ({3.8+2.75*\i},0) {\h};}
\foreach \i/\l/\a/\b/\c/\d in {%
  1/{Simulation and verification}/{in Phase I scope}/{carried forward}/{carried forward}/{same direct work},
  2/{Regulatory and protocol}/{in Phase I scope}/{carried forward}/{carried forward}/{same direct work},
  3/{Site and device readiness}/{not in scope}/{in Phase II scope}/{carried forward}/{same direct work},
  4/{Clinical conduct}/{not in scope}/{partial}/{full}/{same direct work}}{
  \node[d2cellg,minimum width=3.8cm,minimum height=\rh cm,anchor=north west] at (0,{-\rh*\i}) {\l};
  \node[d2cell,minimum width=2.75cm,minimum height=\rh cm,anchor=north west] at (3.8,{-\rh*\i}) {\a};
  \node[d2celll,minimum width=2.75cm,minimum height=\rh cm,anchor=north west] at (6.55,{-\rh*\i}) {\b};
  \node[d2celll,minimum width=2.75cm,minimum height=\rh cm,anchor=north west] at (9.30,{-\rh*\i}) {\c};
  \node[d2cell,minimum width=2.75cm,minimum height=\rh cm,anchor=north west] at (12.05,{-\rh*\i}) {\d};}
\foreach \i/\l/\a/\b/\c/\d in {%
  5/{Award}/{306000 dollars}/{1300000 dollars}/{3500000 dollars}/{2137000 dollars},
  6/{Direct work in the award}/{part of 1396000}/{part of 1396000}/{3500000 dollars}/{1396000 dollars},
  7/{Overhead and fee load}/{15.0 percent}/{15.0 percent}/{direct only}/{57 percent}}{
  \node[d2cellk,minimum width=3.8cm,minimum height=\rh cm,anchor=north west] at (0,{-\rh*\i}) {\l};
  \node[d2celll,minimum width=2.75cm,minimum height=\rh cm,anchor=north west] at (3.8,{-\rh*\i}) {\a};
  \node[d2celll,minimum width=2.75cm,minimum height=\rh cm,anchor=north west] at (6.55,{-\rh*\i}) {\b};
  \node[d2celll,minimum width=2.75cm,minimum height=\rh cm,anchor=north west] at (9.30,{-\rh*\i}) {\c};
  \node[d2cellg2,minimum width=2.75cm,minimum height=\rh cm,anchor=north west] at (12.05,{-\rh*\i}) {\d};}
\draw[trialchar,line width=0.7pt] (0,{-\rh}) -- (14.8,{-\rh});
\draw[trialchar,line width=0.7pt] (0,{-\rh*5}) -- (14.8,{-\rh*5});
\node[d2mid,text width=58mm,anchor=north west] at (0,{-\rh*8-0.20})
  {Federal below a firewall, contributed unpriced, private above: 3.67 to one};
\node[d2soft,text width=52mm,anchor=north west] at (7.60,{-\rh*8-0.20})
  {2.6 FTE at full Phase II staffing, 33 months};
\pnote{0}{-6.85}{Column 5 is the counterfactual: the same direct work routed through an
institution rather than through the company,\\and it is the only column filled
neutral across the money band.}
\end{appfig}
\vspace{-0.6cm}
\figcaption{Figure 18. One four-layer budget priced three ways, and the same
direct work\\under a 15.0 percent company load against a 57 percent university
rate.}
\end{appfloat}

% ---------------------------------------------------------------------------
% FIGURE 19.  graphviz-type, record with ports.  Four port cells at one 0.56 cm
% height, four work package records at 1.80 cm, one firewall rule crossed
% exactly twice.
% ---------------------------------------------------------------------------
\begin{appfloat}[!tb]
\begin{appfig}[graphviz-type / record with ports]
\node[gvcellh,minimum width=4.6cm,minimum height=0.52cm,anchor=north west] at (0,0) {Award, 1606000 dollars over 33 months};
\foreach \i/\p in {0/{Port 1, simulation and verification},1/{Port 2, regulatory and protocol},
  2/{Port 3, site and device readiness},3/{Port 4, clinical conduct, partial}}{
  \node[gvcellk,minimum width=4.6cm,minimum height=0.56cm,anchor=north west]
    at (0,{-0.52-0.56*\i}) {\p};}
\node[gvcellh,minimum width=4.6cm,minimum height=0.52cm,anchor=north west] at (0,-2.76) {direct work inside the award, 1396000};
\foreach \i/\t/\a/\b in {%
  0/{Simulation and verification}/{Phase 0 gate evidence}/{VVUQ suite maintenance},
  1/{Regulatory and protocol}/{IND maintenance}/{amendments and reports},
  2/{Site and device readiness}/{one site qualified}/{device acceptance testing},
  3/{Clinical conduct, partial}/{first cohorts only}/{monitoring and adjudication}}{
  \node[gvcellh,minimum width=4.0cm,minimum height=0.46cm,anchor=north west] at (6.15,{-0.10-1.80*\i}) {\t};
  \node[gvcells,minimum width=4.0cm,minimum height=0.46cm,anchor=north west] at (6.15,{-0.56-1.80*\i}) {\a};
  \node[gvcells,minimum width=4.0cm,minimum height=0.46cm,anchor=north west] at (6.15,{-1.02-1.80*\i}) {\b};}
\foreach \i/\t/\a/\b in {%
  0/{Outside the award}/{contributed drug supply}/{unpriced in the bridge},
  1/{Outside the award}/{operating room and pathology}/{contributed by the site}}{
  \node[gvcellh,minimum width=4.0cm,minimum height=0.46cm,anchor=north west] at (11.35,{-3.70-1.80*\i}) {\t};
  \node[gvcellg,minimum width=4.0cm,minimum height=0.46cm,anchor=north west] at (11.35,{-4.16-1.80*\i}) {\a};
  \node[gvcellg,minimum width=4.0cm,minimum height=0.46cm,anchor=north west] at (11.35,{-4.62-1.80*\i}) {\b};}
\draw[trialchar,dashed,line width=0.8pt] (10.85,0.20) -- (10.85,-7.40);
\node[gvctitle,rotate=90,anchor=south] at (10.75,-3.60) {firewall, federal money left of the line};
\foreach \i in {0,1,2,3}{
  \draw[gvedgeb] (4.60,{-0.80-0.56*\i}) -- (6.15,{-0.56-1.80*\i});}
\draw[gvedged] (10.15,-4.16) -- node[above,font=\tiny,fill=trialwhite,inner sep=1pt] {requires} (11.35,-4.16);
\draw[gvedged] (10.15,-5.96) -- node[above,font=\tiny,fill=trialwhite,inner sep=1pt] {requires} (11.35,-5.96);
\node[gvboxm,text width=54mm,anchor=north west] at (0,-4.05) {2104000 dollars of the five-year program sits outside the award};
\node[gvboxs,text width=40mm,anchor=north west] at (0,-5.35) {5900000 dollars of private capital sits above the firewall};
\pnote{0}{-8.05}{The four port cells share one 0.56 cm height, so the four edges leave at evenly
spaced anchors and their fan is\\regular. Each requirement edge crosses the
firewall exactly once, which is the point of drawing the firewall at all.}
\end{appfig}
\vspace{-0.6cm}
\figcaption{Figure 19. The award as a four-port record, the work package each
port pays,\\and the two packages that sit outside the award and are funded
another way.}
\end{appfloat}

\subsection{The NIH application}

Version 2.0 of the RFA-RM-27-001 application, deposited July 12, 2026, is the
most complete single funding document in the set: twelve section files carrying
the cover and transmittal, opportunity and organization, project summary,
research strategy, clinical trial synopsis, human subjects and regulatory,
statistics and safety and operations, data management and sharing, facilities
and resources, budget and milestones and sustainability, investigator profile,
and assurances and attachments \cite{rfarm27001v2}. Version 1.0 preceded it by
five days \cite{rfarm27001v1}. Table 19 states what each section had to
establish for a reviewer.

\begin{apptable}
\begin{tabularx}{\textwidth}{>{\raggedright\arraybackslash}p{0.85cm} >{\raggedright\arraybackslash}p{5.35cm} >{\raggedright\arraybackslash}X}
\headrow \thead{No.} & \thead{Application section} & \thead{What it had to establish}\\
\midrule
00 & Cover and transmittal & Applicant identity and mechanism fit \\
01 & Opportunity and organization & Why this RFA, and who is applying \\
02 & Project summary & The one-paragraph case \\
03 & Research strategy & Significance, innovation, approach \\
04 & Clinical trial synopsis & The study in the funder's own form fields \\
05 & Human subjects and regulatory & IRB, consent, and the IND status \\
06 & Statistics, safety and operations & Powering, monitoring, and conduct \\
07 & Data management and sharing & Where the data goes and who may read it \\
08 & Facilities and resources & What exists and what must be obtained \\
09 & Budget, milestones and sustainability & The money and the stop conditions \\
10 & Investigator profile & The single investigator's record \\
11 & Assurances and attachments & The certifications the form requires \\
\end{tabularx}
\end{apptable}
\vspace{-0.6cm}
\tabcap{Table 19. The twelve sections of the NIH application at version 2.0, and\\what
each section had to establish for a reviewer.}

% ---------------------------------------------------------------------------
% FIGURE 20.  diagrams (python)-type, clustered infrastructure.  Four aligned
% ranks at one 1.90 cm vertical pitch; converging and diverging fans routed
% through stated clear corridors.
% ---------------------------------------------------------------------------
\begin{appfloat}[!tb]
\begin{appfig}[diagrams (python)-type / clustered infrastructure]
\dgnodeg{dgtileg}{0.85}{1.20}{\glyphdoc}{Protocols and IND}{e1}
\dgnodeg{dgtileg}{0.85}{-0.70}{\glyphcpu}{Simulation record}{e2}
\dgnodeg{dgtileg}{0.85}{-2.60}{\glyphbank}{Legislative drafts}{e3}
\dgnodeg{dgtileg}{0.85}{-4.50}{\glyphteam}{Peer review record}{e4}
\begin{scope}[on background layer]
\node[dgcluster2,fit=(e1)(e1l)(e2)(e2l)(e3)(e3l)(e4)(e4l)] (ce) {};\end{scope}
\node[dgctitle2,anchor=south west] at ([yshift=1.5mm]ce.north west) {Evidence store};
\dgnodew{dgtiled}{5.05}{1.20}{\glyphai}{Master prompt}{p1}
\dgnodew{dgtiled}{5.05}{-0.70}{\glyphgear}{Sub-prompt schedule}{p2}
\dgnodew{dgtiled}{5.05}{-2.60}{\glyphlink}{Draft, full, final stages}{p3}
\dgnodew{dgtiled}{5.05}{-4.50}{\glyphshield}{Defect and format pass}{p4}
\begin{scope}[on background layer]
\node[dgcluster,fit=(p1)(p1l)(p2)(p2l)(p3)(p3l)(p4)(p4l)] (cp) {};\end{scope}
\node[dgctitle,anchor=south west] at ([yshift=1.5mm]cp.north west) {Pipeline};
\dgnode{dgtilem}{9.85}{1.20}{\glyphuser}{Person-based, 2 applications}{x1}
\dgnode{dgtilem}{9.85}{-0.70}{\glyphserver}{Organization, 3 applications}{x2}
\dgnode{dgtilem}{9.85}{-2.60}{\glyphhand}{Partnership, 4 applications}{x3}
\dgnode{dgtilem}{9.85}{-4.50}{\glyphflask}{Small business, 1 plus the plan}{x4}
\begin{scope}[on background layer]
\node[dgcluster,fit=(x1)(x1l)(x2)(x2l)(x3)(x3l)(x4)(x4l)] (cx) {};\end{scope}
\node[dgctitle,anchor=south west] at ([yshift=1.5mm]cx.north west) {Typed exits};
\dgnode{dgtile}{14.10}{0.25}{\glyphlink}{Zenodo DOI}{d1}
\dgnode{dgtile}{14.10}{-1.65}{\glyphdb}{Repository path}{d2}
\dgnode{dgtile}{14.10}{-3.55}{\glyphlock}{Commit provenance}{d3}
\begin{scope}[on background layer]
\node[dgcluster,fit=(d1)(d1l)(d2)(d2l)(d3)(d3l)] (cd) {};\end{scope}
\node[dgctitle,anchor=south west] at ([yshift=1.5mm]cd.north west) {Deposit};
\foreach \n in {e1,e2,e3,e4}{\draw[dgedge] (\n.east) -- (p1.west);}
\draw[dgedgeb] (p1l.south) -- (p2.north);
\draw[dgedgeb] (p2l.south) -- (p3.north);
\draw[dgedgeb] (p3l.south) -- (p4.north);
\foreach \n in {x1,x2,x3,x4}{\draw[dgedgeb] (p4.east) -- (\n.west);}
\foreach \n in {x1,x2,x3,x4}{\draw[dgedged] (\n.east) -- (cd.west);}
\pnote{0.35}{-6.30}{The four inbound edges converge through the clear corridor between x = 2.35 and
x = 3.55 and the four outbound\\edges diverge through x = 6.55 to 8.35; the four
deposit edges meet one stated waypoint outside the deposit cluster.}
\end{appfig}
\vspace{-0.6cm}
\figcaption{Figure 20. One evidence store, one production pipeline, and the four
typed\\exits at which a funding artifact leaves it for a named recipient class.}
\end{appfloat}

\subsection{What the same work costs each way}

The cost case is not an argument about effort; it is a comparison of published
benchmarks against a projected figure the author's own simulation record
supports. An empirical virtual trial in triplicate is benchmarked above
\$120,000 per run and about 4.5 months; a ten-arm quantitative systems
pharmacology virtual trial above \$2,000,000 and several months to over a year;
a digital-twin simulator platform between \$28,000 and \$700,000 and 3 to 16
weeks \cite{capplan2026, qspmetpancre, digitaltwinfda2025}. Against all three,
the author's comparable work is projected at \$36,330 and about one month. The
word is projected rather than estimated, because the figure is a projection from
the source set rather than a measured invoice, and the capitalization plan
reports it that way \cite{capplan2026}. Table 20 sets the five activities side
by side.

\begin{apptable}
\begin{tabularx}{\textwidth}{>{\raggedright\arraybackslash}p{4.35cm} >{\raggedright\arraybackslash}p{3.15cm} >{\raggedright\arraybackslash}p{2.55cm} >{\raggedright\arraybackslash}X}
\headrow \thead{Activity} & \thead{Industry benchmark} & \thead{Industry time} & \thead{This work}\\
\midrule
Empirical virtual trial, triplicate & Above \$120,000 per run & 4.5 months & \$36,330 projected, about one month \\
Ten-arm QSP virtual trial & Above \$2,000,000 & Several months to over a year & \$36,330 projected, about one month \\
Digital-twin simulator platform & \$28,000 to \$700,000 & 3 to 16 weeks & \$36,330 projected, about one month \\
Phase 1 IND dossier & Group effort & Months & 4 days, one author \\
Phase 1 and Phase 2 protocols & Group effort & 6 to 12 months each & 2 days each \\
\end{tabularx}
\end{apptable}
\vspace{-0.6cm}
\tabcap{Table 20. Five activities priced against published industry benchmarks,\\with
the projected cost and elapsed time this work reached on each.}

\subsection{What the applications concede}

A funding section that reports only favorable numbers is not useful to a funder,
and the capitalization plan does not read that way. It names four concessions in
its own risk section, and they are worth restating because a reader of this
paper will reach them anyway. The individual-scientist framing that made all ten
prior applications distinctive is abandoned when the applicant becomes a
company. Capitalization detail the author might prefer to keep private is
published permanently under a digital object identifier. Three of the ten prior
recipients will read a company plan as off-mechanism. And the plan proposes
\$1,606,000 and 2.6 full-time equivalents against a novel-performers chapter that
describes work at tens of millions of dollars with teams of ten to a hundred
people over half a decade \cite{capplan2026, kratsios2026sciencegoldenage}.

The last of those four is the one that decides the mechanism. Figure 20's typed
exits are not a taxonomy for its own sake: nine of the ten applications address
mechanisms that fund a person or an institution, and a firm applying to a
person-based mechanism is off-mechanism however good the science is. That is why
application 05 and the capitalization plan that rewrites it are the only two
artifacts in the set aimed at the exit a company can actually reach.

\subsection{What a funder can verify before committing}

Every quantity in this section is checkable before any money moves, which is the
property the deposit discipline of \S2 exists to produce. The ten applications
are deposited with their mechanism class and their ask
\cite{tenapps2026}. The capitalization plan is deposited with its four-layer
budget, its two overhead regimes, its twelve costed milestones and its stop
conditions \cite{capplan2026}. Both NIH application versions are deposited five
days apart with their section structure intact
\cite{rfarm27001v1, rfarm27001v2}. The simulation record behind the \$36,330
projection is deposited, and so is the QSP work it derives from
\cite{qspmetpancre, digitaltwinfda2025}.

A funder who wants to test the production-rate claim rather than accept it can
do so in three steps: read Figure 17's calendar, open the commit history of the
directory named on each row, and compare the dates. A funder who wants to test
the cost claim can read the benchmark column of Table 20 against the cited
sources and then read the author's own simulation record. Neither test requires
the author's cooperation, and neither is available for a plan that exists only
as a narrative.
